\documentclass[11pt]{article}
\usepackage[a4paper,margin=26mm]{geometry}
\usepackage{amsmath,amssymb}
\usepackage[hidelinks]{hyperref}
\title{Conjecture 00000000821: no uniform gradient lower bound at the Laplace eigenvalue scale}
\author{lizaixi01, with OpenAI Codex assistance}
\date{22 October 2026}
\begin{document}
\maketitle
\section{Source and the spectral convention}
The source conjectures a positive uniform lower bound for
\[
 \frac{\max |\nabla\varphi_\lambda|}
 {\lambda\max|\varphi_\lambda|}.
\]
It does not specify the underlying operator or its spectral convention.
We use the standard positive Laplacian convention
$-\Delta\varphi_\lambda=\lambda\varphi_\lambda$, with
$\lambda$ the eigenvalue itself. This is, for example, the convention used
in the abstract of Steinerberger's paper \cite{steinerberger}.
Our conclusion concerns this interpretation. If $\lambda$ instead means
the square root of the eigenvalue, the statement is a different claim
and the present argument does not refute it.

The counterexamples lie on one fixed compact manifold: the flat circle
$M=\mathbb R/(2\pi\mathbb Z)$ with arc-length metric $d\theta^2$.
Thus even allowing the lower-bound constant to depend on this manifold
does not rescue the assertion. In these coordinates the gradient norm is
$|f'(\theta)|$ and the positive Laplacian is $-f''(\theta)$.
These formulas follow directly from the constant metric coefficient
$g_{11}=g^{11}=1$. The usual sine/cosine spectrum of this circle is also
recorded in Tsogtgerel's lecture notes \cite{tsogtgerel}, with the opposite
sign convention for the Laplacian.

\section{A family on the actual circle}
For each integer $m\geq1$ define
\[
 f_m([\theta])=\sin(m\theta),\qquad
 D_m([\theta])=m\cos(m\theta).
\]
Both expressions are $2\pi$-periodic, hence well-defined on $M$.
They are continuous, their coordinate pullbacks are smooth, and
\[
 f_m'=D_m,\qquad D_m'=-m^2f_m.
\]
Consequently $f_m$ is a nonzero smooth eigenfunction of the positive
Laplacian, with eigenvalue $\lambda_m=m^2>0$.

The elementary bounds $|\sin t|\leq1$ and $|\cos t|\leq1$, together with
the points $[\pi/(2m)]$ and $[0]$, give the attained maxima
\[
 \max_M|f_m|=1,\qquad \max_M|\nabla f_m|=m.
\]
No supremum is substituted for an unattained maximum. Therefore
\[
 \frac{\max_M|\nabla f_m|}{\lambda_m\max_M|f_m|}
 =\frac{m}{m^2}=\frac1m.                         \tag{1}
\]
For every $c>0$, the Archimedean property supplies $m\geq1$ with $1/m<c$.
Equation (1) rules out every positive constant uniform over eigenfunctions,
including one uniform only on the high-frequency tail. Amplitude
normalization has no effect on this ratio. The exhibited functions already
have absolute maximum one, so the denominator-placement ambiguity in the
Chinese text has no effect on this family.

\section{What is formalized}
The Lean object \texttt{FlatCircle} is Mathlib's actual additive quotient
\texttt{AddCircle (2 * Real.pi)}. The functions \texttt{circleU} and
\texttt{circleD} are defined by the periodic-function quotient lift.
Their pullbacks are proved equal to the sine and cosine expressions;
quotient-map continuity and smoothness of the first pullback are proved.

\texttt{CircleEigenpair} states positivity of the eigenvalue, continuity
of both circle functions, a nonzero value, smoothness of the pullback,
and two derivative identities in Lean's \texttt{HasDerivAt} predicate.
It is the explicit arc-length coordinate formulation described above.
We do not claim that the code invokes a general Riemannian gradient or
Laplace--Beltrami API. The metric-to-coordinate formulas are spelled out
here. No spectral, derivative, or maximum identity is assumed
in the proof of the example.

Theorems \texttt{max\_value} and \texttt{max\_gradient} use
\texttt{IsGreatest} on the value ranges of the actual quotient functions.
The final theorem \texttt{no\_uniform\_positive\_bound} negates the fully
quantified positive bound for these nonzero circle eigenfunctions. Its
proof constructs an arbitrary-small-ratio family, rather than using the
zero function or a single finite-frequency discrepancy.

\section{Verification and provenance}
Lean 4.33.0 and Mathlib are pinned in the accompanying project. Full
project build, source-bound axiom audit and official fresh-kernel replay
are recorded in the evidence directory. Only the usual Lean foundational
axioms are permitted; no admitted proof, native decision shortcut, or new
axiom is used. The PDF was rendered and every page inspected.

The bilingual source is preserved from upstream commit:\par\noindent
\texttt{b355e42906e6eaf52bc7e9d005e0ef02d8aa7400}.
OpenAI Codex assisted the mathematics, formalization and report. A separate
supervising Codex agent reviewed the source and interpretation; this is AI
review, not human certification. The family and its spectral behavior are
standard; we make no claim of new mathematics. Organizer acceptance remains
pending.

\begin{thebibliography}{2}
\bibitem{steinerberger} S. Steinerberger,
\emph{On the Spectral Resolution of Products of Laplacian Eigenfunctions},
2017. \url{https://arxiv.org/abs/1711.09826}.
\bibitem{tsogtgerel} G. Tsogtgerel, \emph{Math 319, Lecture 25},
McGill University, 8 March 2011, slide 2.
\url{https://www.math.mcgill.ca/gantumur/math319w11/downloads/slides25.pdf}.
\end{thebibliography}
\end{document}
